% P15-U1: direct weak-law proofs, all initial applications and exact tail criterion.
\chapter{Weak Laws and Triangular Arrays}
\label{sec:p15-weak-laws-and-triangular-arrays}

Weak laws describe concentration in probability. Their direct variance
and truncation arguments remain useful when rows have different laws,
when observations have no second moment, and when the deterministic
center changes with sample size. This chapter retains those arguments
and the source's polynomial, geometric and combinatorial applications.

The source is Durrett's author Version 5 dated 11 January 2019,
Section 2.2 through Theorem 2.2.12 and its continued proof, with
the immediate integrable corollary 2.2.14
\cite[pp.~56--64]{DurrettPTEAuthor2019}.
The necessity direction of the tail criterion is proved here by
independent copies and a complete symmetric maximal inequality.
The separate historical reading of Feller, Volume II, second edition
(1971), pp.~234--236, cited by the source for that direction, is
recorded at the chapter's boundary.

\section{The direct variance method}
\label{sec:p15-direct-weak-variance}

\begin{hbdefinition}{def:p15-weak-convergence-probability}{Convergence in probability and rows of observations}
For real measurable \(Z_n,Z\) on a common probability space,
\(Z_n\to Z\) in probability means
\(\mathbb P(|Z_n-Z|>\varepsilon)\to0\) for every \(\varepsilon>0\).
Convergence in mean square means
\(\mathbb E|Z_n-Z|^2\to0\).
When the target is a constant, the deviation-probability formulation
also makes sense if the variables for different \(n\) live on different
probability spaces.

A triangular array has real entries \(X_{n,k}\), \(1\leq k\leq n\).
Row independence means that each finite family
\(X_{n,1},\ldots,X_{n,n}\) is mutually independent. Independence
between rows is not required. Variables with finite second moments
are uncorrelated when
\(\mathbb E(X_iX_j)=\mathbb EX_i\,\mathbb EX_j\) for \(i\ne j\).
\index{weak law of large numbers}\index{triangular array}
\end{hbdefinition}

\begin{hbtheorem}{thm:p15-direct-variance-weak-laws}{proved=annex:p15.direct-variance-weak-laws}{The four direct variance and moment statements}
\begin{enumerate}
\item For uncorrelated finite-second-moment \(X_1,\ldots,X_n\),
\(\operatorname{Var}(\sum_iX_i)=\sum_i\operatorname{Var}(X_i)\).
\item If \(p>0\) and \(\mathbb E|Z_n|^p\to0\), then
\(Z_n\to0\) in probability; explicitly, for each \(\varepsilon>0\),
\(\mathbb P(|Z_n|\geq\varepsilon)\leq
\varepsilon^{-p}\mathbb E|Z_n|^p\).
\item If \(X_i\) are uncorrelated, have common finite mean \(\mu\),
and \(\operatorname{Var}(X_i)\leq C<\infty\), then, for
\(S_n=\sum_{i=1}^nX_i\),
\[
 \mathbb E[(S_n/n-\mu)^2]\leq C/n,
 \qquad S_n/n\longrightarrow\mu
       \quad\text{in mean square and in probability}.
\]
\item For any real \(S_n\) with finite second moment, let
\(\mu_n=\mathbb ES_n\), \(\sigma_n^2=\operatorname{Var}(S_n)\),
and let \(b_n\ne0\) be finite real constants. If
\(\sigma_n^2/b_n^2\to0\), then
\((S_n-\mu_n)/b_n\to0\) in mean square and in probability.
This last assertion requires no representation as a sum and no
independence assumption.
\end{enumerate}
These retain source Theorem 2.2.1, Lemma 2.2.2 and Theorems 2.2.3
and 2.2.6 with their direct proofs.
\end{hbtheorem}

\begin{hbtheorem}{thm:p15-bernstein-approximation}{proved=annex:p15.bernstein-approximation}{The Bernstein polynomial application}
Let \(f:[0,1]\to\mathbb R\) be continuous and define, for \(n\geq1\),
\[
 f_n(x)=\sum_{m=0}^n\binom nm x^m(1-x)^{n-m}f(m/n),
 \qquad \binom nm=\frac{n!}{m!(n-m)!}.
\]
Powers with exponent zero are the constant polynomial one, including
at zero, and \(0!=1\). These polynomials have degree at most \(n\),
and converge uniformly to \(f\).
More precisely, if \(M_f=\sup_{[0,1]}|f|\) and
\(\omega_f(\delta)=\sup_{|u-v|\leq\delta}|f(u)-f(v)|\), with
\(u,v\in[0,1]\), then
\[
 \sup_{x\in[0,1]}|f_n(x)-f(x)|
       \leq\omega_f(\delta)+\frac{M_f}{2n\delta^2}
                        \quad(\delta>0).
\]
For \(f(x)=x^2\), the polynomial is
\(f_n(x)=x^2+x(1-x)/n\), and its exact uniform error is \(1/(4n)\).
\end{hbtheorem}

\begin{hbtheorem}{thm:p15-cube-shell-concentration}{proved=annex:p15.cube-shell-concentration}{The high-dimensional cube application}
Let \(\lambda_n\) be product Lebesgue measure on \(\mathbb R^n\),
let \(\|x\|_2=(\sum_i x_i^2)^{1/2}\), and set
\[
 A_{n,\varepsilon}=
 \{x:(1-\varepsilon)\sqrt{n/3}<\|x\|_2
                              <(1+\varepsilon)\sqrt{n/3}\}.
\]
For every \(\varepsilon>0\),
\[
 \frac{\lambda_n(A_{n,\varepsilon}\cap(-1,1)^n)}{2^n}
       \longrightarrow1.
\]
The portion of the cube outside this annulus has relative volume at most
\(\min\{1,4/(5n\varepsilon^2)\}\).
The statement concerns normalized volume in an annulus, not equality
of the cube with a sphere.
\end{hbtheorem}

\section{Three applications to changing rows}
\label{sec:p15-changing-row-applications}

\begin{hbdefinition}{def:p15-collector-and-cycle-models}{Coupon times and canonical cycle order}
For each positive integer \(n\), coupon draws are i.i.d. uniform on
\(\{1,\ldots,n\}\). For \(1\leq k\leq n\), let \(\tau_k^n\) be the first draw index at
which \(k\) different coupons have appeared, set \(\tau_0^n=0\),
and assign an unattained hitting time \(+\infty\). The proof establishes
\(\tau_n^n<\infty\) almost surely. On that event set
\(D_{n,k}=\tau_k^n-\tau_{k-1}^n\); on its null complement set
every \(D_{n,k}=1\). Put \(T_n=\sum_{k=1}^nD_{n,k}\), a finite
version of the collection time agreeing with \(\tau_n^n\) almost surely.

For a uniform permutation of \(\{1,\ldots,n\}\), list each cycle
starting at its smallest unused label, follow the permutation until
that cycle closes, and begin the next cycle at the smallest unused
label. Let \(I_{n,k}\) indicate a closing parenthesis after the
\(k\)-th label in this concatenated list, and let
\(K_n=\sum_{k=1}^n I_{n,k}\) be the number of cycles.
Write \(H_n=\sum_{j=1}^n1/j\).
A geometric law with success probability \(p\in(0,1]\) is the
law on positive integers with masses \(p(1-p)^{d-1}\), \(d\geq1\);
at \(p=1\) it is interpreted as point mass at one.
\end{hbdefinition}

\begin{hbtheorem}{thm:p15-coupon-collector-law}{proved=annex:p15.coupon-collector-law}{Coupon waiting times and their weak law}
The durations \(D_{n,1},\ldots,D_{n,n}\) are independent, with
geometric laws on \(\{1,2,\ldots\}\) and success probabilities
\(p_{n,k}=(n-k+1)/n\). The case \(p_{n,1}=1\) means
\(D_{n,1}=1\) almost surely. Moreover,
\[
 \mathbb ET_n=nH_n,\qquad
 \operatorname{Var}(T_n)=n^2\sum_{j=1}^n j^{-2}-nH_n\leq2n^2.
\]
As \(n\to\infty\), with the denominator used for \(n\geq2\),
\[
 \frac{T_n-nH_n}{n\log n}\longrightarrow0,\qquad
 \frac{T_n}{n\log n}\longrightarrow1
                                     \quad\text{in probability}.
\]
For \(n=365\), the leading scale \(365\log365\) is approximately
\(2153.46\), while the exact finite expectation \(365H_{365}\) is
approximately \(2364.65\). These are distinct quantities in the
specified uniform-coupon model.
\end{hbtheorem}

\begin{hbtheorem}{thm:p15-permutation-cycle-law}{proved=annex:p15.permutation-cycle-law}{Independent closing indicators and cycle-count concentration}
The indicators \(I_{n,1},\ldots,I_{n,n}\) are independent and
\(\mathbb P(I_{n,k}=1)=1/(n-k+1)\). Thus
\[
 \mathbb EK_n=H_n,\qquad
 v_n=\operatorname{Var}(K_n)=H_n-\sum_{j=1}^n j^{-2}\leq H_n.
\]
For every \(\varepsilon>0\),
\[
 \frac{K_n-H_n}{(\log n)^{1/2+\varepsilon}}\longrightarrow0,
 \qquad K_n/\log n\longrightarrow1
                                     \quad\text{in probability}.
\]
The logarithmic denominators start at \(n=2\). At exponent \(1/2\)
the centered conclusion fails: for all sufficiently large \(n\),
\[
 \mathbb P(|K_n-H_n|/\sqrt{\log n}>1/2)\geq1/16.
\]
For the source permutation
\(\pi(1),\ldots,\pi(9)=(3,9,6,8,2,1,5,4,7)\),
the canonical cycles are \((1\,3\,6)(2\,9\,7\,5)(4\,8)\);
the closing positions are \(3,7,9\).
\end{hbtheorem}

\begin{hbtheorem}{thm:p15-occupancy-weak-law}{proved=annex:p15.occupancy-weak-law}{Empty boxes and dependent indicators}
Place a nonnegative integer number \(r=r_n\) of labelled balls independently and uniformly into
\(n\) boxes, so all \(n^r\) assignments are equally likely.
Let \(N_n\) count empty boxes and, for \(n\geq2\), set
\(q_n=(1-1/n)^r\), \(s_n=(1-2/n)^r\).
When \(r=0\), these powers are empty products and equal one.
Then
\[
 \mathbb EN_n=nq_n,\qquad
 \operatorname{Var}(N_n)
    =nq_n(1-q_n)+n(n-1)(s_n-q_n^2)\leq n/4.
\]
If \(r_n/n\to c\in[0,\infty]\), then
\(N_n/n\to e^{-c}\) in probability, with \(e^{-\infty}=0\).
The finite-\(c\) assertion is the source application; the infinite
case follows from the same bounds. In particular \(r_n=n\)
gives limiting empty-box proportion \(e^{-1}\).
Empty-box indicators are generally dependent.
\end{hbtheorem}

\section{Truncation and the full tail criterion}
\label{sec:p15-truncation-weak-criterion}

\begin{hbdefinition}{def:p15-array-truncation}{The source row truncation and its center}
For a row-independent triangular array and finite positive \(b_n\),
put
\[
 \overline X_{n,k}=X_{n,k}\mathbf1_{\{|X_{n,k}|\leq b_n\}},
 \quad S_n=\sum_{k=1}^nX_{n,k},\quad
 \overline S_n=\sum_{k=1}^n\overline X_{n,k},\quad
 a_n=\sum_{k=1}^n\mathbb E\overline X_{n,k}.
\]
The truncation deletes values beyond the cutoff and retains values
exactly at it. The center uses truncated expectations, which exist
even if the original row entries are not integrable.
\end{hbdefinition}

\begin{hbtheorem}{thm:p15-triangular-array-weak-law}{proved=annex:p15.triangular-array-weak-law}{The source weak law for triangular arrays}
Assume \(b_n>0\), \(b_n\to\infty\), row independence, and
\[
 u_n=\sum_{k=1}^n\mathbb P(|X_{n,k}|>b_n)\longrightarrow0,\qquad
 w_n=b_n^{-2}\sum_{k=1}^n\mathbb E\overline X_{n,k}^{\,2}
                                                 \longrightarrow0.
\]
Then \((S_n-a_n)/b_n\to0\) in probability. For each
\(\varepsilon>0\), the direct bound is
\[
 \mathbb P(|S_n-a_n|>\varepsilon b_n)
                            \leq u_n+\varepsilon^{-2}w_n.
\]
There is no independence requirement between rows. These two
truncation conditions are sufficient tests, not an if-and-only-if
characterization of every array with a weak law.
\end{hbtheorem}

\begin{hblemma}{thm:p15-levy-symmetric-maximal}{proved=annex:p15.levy-symmetric-maximal}{The symmetric maximal inequality}
Let \(Y_1,\ldots,Y_n\) be independent real random variables whose
individual laws are symmetric about zero, and put
\(T_0=0\), \(T_k=\sum_{i=1}^kY_i\). For every \(t>0\),
\[
 \mathbb P\!\left(\max_{1\leq k\leq n}|T_k|>t\right)
          \leq2\mathbb P(|T_n|>t),\qquad
 \mathbb P\!\left(\max_{1\leq k\leq n}|Y_k|>2t\right)
          \leq2\mathbb P(|T_n|>t).
\]
No moment or continuity assumption is needed. The first inequality
also holds with both strict inequalities replaced by non-strict ones.
\end{hblemma}

\begin{hbtheorem}{thm:p15-iid-weak-tail-criterion}{proved=annex:p15.iid-weak-tail-criterion}{The full weak-law criterion and all deterministic centers}
Let \(X_n\) be i.i.d. real random variables, set \(S_n=\sum_{i=1}^nX_i\),
and define the finite truncated means
\(\mu_n=\mathbb E[X_1\mathbf1_{\{|X_1|\leq n\}}]\).
There exist real deterministic \(c_n\) with
\(S_n/n-c_n\to0\) in probability if and only if
\[
 x\mathbb P(|X_1|>x)\longrightarrow0\qquad(x\to\infty).
\]
Under this condition the source center always works:
\(S_n/n-\mu_n\to0\) in probability. More precisely, for an
arbitrary deterministic sequence \(c_n\),
\[
 S_n/n-c_n\longrightarrow0\ \text{in probability}
 \quad\Longleftrightarrow\quad
 \begin{cases}
 x\mathbb P(|X_1|>x)\to0,\\
 c_n-\mu_n\to0.
 \end{cases}
\]
The tail condition implies
\(\mathbb E|X_1|^r<\infty\) for every \(0<r<1\);
it need not imply a finite first absolute moment.
Sufficiency is the source's direct triangular-array proof.
Necessity is proved independently through the symmetric maximal
inequality, with no central limit theorem.
\end{hbtheorem}

\begin{hbcorollary}{thm:p15-direct-integrable-weak-law}{proved=annex:p15.direct-integrable-weak-law}{The integrable weak law by the source truncation proof}
If \(X_n\) are i.i.d. and \(\mathbb E|X_1|<\infty\), then
\[
 \frac1n\sum_{i=1}^nX_i\longrightarrow\mathbb EX_1
                                         \quad\text{in probability}.
\]
Its proof uses two dominated-convergence calculations and the
tail criterion, retaining source Theorem 2.2.14's separate argument.
\end{hbcorollary}

\section{What the hypotheses do and do not imply}
\label{sec:p15-weak-law-boundaries}

\begin{hbtheorem}{thm:p15-weak-law-counterexamples}{proved=annex:p15.weak-law-counterexamples}{Truncation failures, cancellation and a weak law without a mean}
The following array examples use \(b_n=n\).
\begin{enumerate}
\item If \(X_{n,1}=2n\) and every other entry is zero, then
\(w_n=0\), but \(u_n=1\) and \((S_n-a_n)/b_n=2\).
\item If \(X_{n,1}=n\varepsilon\), where \(\varepsilon\) is a
fair sign, and other entries are zero, then the row is independent
and \(u_n=0\), but \(w_n=1\) and the normalized centered sum
is \(\varepsilon\). Equality at the cutoff is retained.
\item If every entry in row \(n\) is the same fair sign, then
\(u_n=0\) and \(w_n=1/n\), but the normalized centered sum is
again that sign. Row independence has failed.
\item Conversely, the independent deterministic row with entries
\(2n,-2n,0,\ldots,0\), for \(n\geq2\), has \(u_n=2\)
but \(S_n=a_n=0\). Thus failing the tail test does not by itself
prove failure of a weak law.
\item For the i.i.d. symmetric law giving mass
\(1/(2\log2\,j2^j)\) to each of \(2^j\) and \(-2^j\),
\(j\geq1\), both part expectations are infinite, yet
\[
 S_n/n\longrightarrow0\quad\text{in probability},\qquad
 \limsup_n|S_n|/n=\infty\quad\text{almost surely}.
\]
\item For an i.i.d. sequence with the positive dyadic law
\(\mathbb P(X_1=2^j)=2^{-j}\),
no deterministic sequence centers \(S_n/n\) to zero in probability.
Indeed \(x\mathbb P(X_1>x)\) equals one at \(x=2^j\).
\end{enumerate}
\end{hbtheorem}

\begin{proofboundary}{def:p15-weak-laws-next-boundary}{The separate historical source and later distributional proofs}
The direct variance, combinatorial, triangular-array and weak-tail
proofs above are complete. The independently supplied necessity proof
establishes the criterion independently. Feller's separate
median-based symmetrization route and all its cited inequalities are
supplied in \cref{thm:p15-feller-symmetrization-inequalities,
thm:p15-feller-largest-term-inequality,thm:p15-feller-generalized-weak-law}.
This is the Volume II, second-edition presentation on pp.~234--236
identified by Durrett's remark \cite[p.~63]{DurrettPTEAuthor2019},
with its own complete source argument
\cite[ch.~VII, sec.~7]{Feller1971ProbabilityII}.

The Cauchy and St.~Petersburg source applications, Examples~2.2.15--16,
are supplied in \cref{thm:p15-cauchy-weak-law-failure,
thm:p15-petersburg-array-calculation,thm:p15-petersburg-normalizers-and-prices}.
The distinct de Moivre--Laplace route is
\cref{thm:p15-stirling-de-moivre-laplace}; the weak-convergence
presentation is \cref{sec:p15-weak-convergence-of-laws}.
Characteristic-function inversion and continuity are
\cref{thm:p15-characteristic-inversion-uniqueness,
thm:p15-characteristic-continuity-theorem}. The remaining
characteristic-function criteria, moment problems and section exercises
are supplied in \cref{sec:p15-characteristic-function-criteria-and-moments}.
The i.i.d. and array central limit proofs, together with the source exercises,
are \cref{sec:p15-central-limit-theorems}. The normal-domain
characterization, including its necessity direction, is
\cref{thm:p15-nda-characterization}; the original truncated-mean
cutoff route is \cref{thm:p15-nda-gk-cutoff-route}.
The cycle-count endpoint failure above was proved by a fourth-moment
argument; the source's distinct central-limit treatment in Example~3.4.11
is supplied in \cref{thm:p15-clt-records-and-cycles}.
\end{proofboundary}
